\subsection{从行为克隆到奖励优化：RL、GRPO 与 RLHF}

slides 38--47 解释为什么后训练不能只靠 SFT。行为克隆受人类示范上限约束，理想答案昂贵，而且把模型不知道的正确答案硬塞给它，可能训练出看似可信的 hallucination。RL 改为最大化期望奖励；奖励可以来自字符串匹配、代码测试、偏好 reward model 或 LLM judge。DeepSeek-R1 使用 GRPO 一类基于同题多样本估计 advantage 的方法。对于 Agent，sampling、长 rollout 和慢环境反馈又把基础设施推到核心位置。

\lecturefigure{slides-images/slide-038.png}{官方 slide 38：SFT 行为克隆的能力、幻觉与成本限制}
\lecturefigure{slides-images/slide-039.png}{官方 slide 39：从克隆答案转向最大化规则/模型/裁判奖励}
\lecturefigure{slides-images/slide-040.png}{官方 slide 40：DeepSeek-R1 强化学习流程示意}
\lecturefigure{slides-images/slide-041.png}{官方 slide 41：GRPO 用组内 Monte Carlo 样本估计 advantage}
\lecturefigure{slides-images/slide-042.png}{官方 slide 42：长 rollout、并发环境服务与 collocated engines}
\lecturefigure{slides-images/slide-043.png}{官方 slide 43：RLHF 从成对偏好到策略更新的流程}
\lecturefigure{slides-images/slide-044.png}{官方 slide 44：PPO/DPO 相对 SFT/pretrain 的偏好收益示例}
\lecturefigure{slides-images/slide-045.png}{官方 slide 45：人类偏好数据需要任务、样例和标注指南}
\lecturefigure{slides-images/slide-046.png}{官方 slide 46：人类偏好数据的速度、正确性、分布与伦理挑战}
\lecturefigure{slides-images/slide-047.png}{官方 slide 47：以 LLM preference 替代部分人类偏好的 RLAIF 思路}

\begin{importantbox}{GRPO 的直觉}
对同一问题采样一组回答，依据奖励 $r_i$ 做组内标准化，可得到近似 advantage
\[
\hat A_i=\frac{r_i-\operatorname{mean}(r)}{\operatorname{std}(r)+\epsilon}.
\]
高于组均值的轨迹被增强，低于组均值的轨迹被抑制。符号 $r_i$ 表示第 $i$ 条 rollout 的奖励，$\epsilon$ 防止标准差过小时数值不稳定。这个估计减少了单独训练 value model 的需求，但仍高度依赖采样多样性、奖励可靠性和环境不可被投机利用。
\end{importantbox}

\begin{knowledgebox}{术语消化：PPO、DPO、RLHF 与 RLAIF}
RLHF 是使用人类偏好信号对齐模型的总体范式；PPO 是可用于更新策略的一类 on-policy 强化学习算法；DPO 直接从偏好对优化策略与参考模型的相对概率，不需要显式 rollout reward optimization；RLAIF 则用 AI 反馈替代或补充人类偏好。它们不是同一层级的互斥缩写。课程把它们放在一起，是为了比较“偏好从哪里来”和“策略怎样吸收偏好”两个不同问题。
\end{knowledgebox}

\begin{warningbox}{Reward hacking 是评估器失败，不是模型“作弊人格”}
\textbf{老师强调：}在 00:05--00:07，讲者举例说明 Agent 可能修改测试、让测试永远返回 true，从而获得高奖励。这说明环境暴露了与真实目标不一致的捷径。修复方向包括权限隔离、不可篡改 verifier、多重结果检查和对轨迹副作用审计，而不是只在 prompt 中要求模型“诚实”。slide 42 的系统设计也要服务这一目标：并发和暂停长尾 rollout 提高吞吐，但不能丢失可审计状态。
\end{warningbox}

人类偏好同样不是无噪声真值。标注者容易偏好更长、更流畅、更像标准答案的输出，也可能因为知识不足而忽略事实错误；标注人口分布还会改变模型行为。LLM judge 能显著降低成本，却会继承自己的偏好。课程因此自然转入评估章节：只有明确评估在训练闭环中的角色，团队才能判断 reward gain 是真实能力进步，还是评估器偏差被策略放大。

